\documentclass[english,hebrew]{article}
\usepackage[HE8,T1]{fontenc}
\usepackage[utf8]{inputenc}
\setcounter{secnumdepth}{-2}
\setlength{\parindent}{0bp}
\usepackage{amsmath}
\usepackage{amsthm}
\usepackage{amssymb}
\PassOptionsToPackage{normalem}{ulem}
\usepackage{ulem}

\makeatletter
\theoremstyle{plain}
\newtheorem{thm}{\protect\theoremname}
\theoremstyle{definition}
\newtheorem{xca}[thm]{\protect\exercisename}
\theoremstyle{definition}
\newtheorem*{sol*}{\protect\solutionname}

\@ifundefined{date}{}{\date{}}
\usepackage{graphicx}
\usepackage{geometry}
\newcommand{\limi}{\underset{n\rightarrow\infty}{\lim}}
\newcommand{\limxz}{\underset{x\rightarrow x_{0}}{\lim}}
\newcommand{\limfxz}{\underset{x\rightarrow x_{0}}{\lim}f(x)}
\newcommand{\limfxm}{\underset{x\rightarrow x_{0}^{-}}{\lim}f(x)}
\newcommand{\limfxp}{\underset{x\rightarrow x_{0}^{+}}{\lim}f(x)}
\newcommand{\limxm}{\underset{x\rightarrow x_{0}^{-}}{\lim}}
\newcommand{\limxp}{\underset{x\rightarrow x_{0}^{+}}{\lim}}
\newcommand{\sqrm}{M_{m\times m}(\mathbb{F})}
\newcommand{\seqa}{(a_{n})_{n=1}^{\infty}}
\newcommand{\seqx}{(x_{n})_{n=1}^{\infty}}
\newcommand{\funcdr}{f:D\rightarrow\mathbb{R}}
\newcommand{\der}{\limxz\frac{f(x)-f(x_{0})}{x-x_{0}}}
\usepackage[all]{pdfprivacy}

\makeatother

\usepackage{babel}
\addto\captionsenglish{\renewcommand{\exercisename}{Exercise}}
\addto\captionsenglish{\renewcommand{\solutionname}{Solution}}
\addto\captionsenglish{\renewcommand{\theoremname}{Theorem}}
\addto\captionshebrew{\renewcommand{\exercisename}{תרגיל}}
\addto\captionshebrew{\renewcommand{\solutionname}{פתרון}}
\addto\captionshebrew{\renewcommand{\theoremname}{משפט}}
\providecommand{\exercisename}{תרגיל}
\providecommand{\solutionname}{פתרון}
\providecommand{\theoremname}{משפט}

\begin{document}
\newgeometry{top=1cm, bottom=2cm, left=1cm, right=1cm}
\title{לקט שאלות - האינטגרל המסוים}
\author{נעם לב}
\maketitle
\begin{abstract}
חלק מהשאלות )ציינתי אילו( נכתבו על ידי צוות הקורס באוניברסיטה העברית.
הפתרון הוא שלי ולא רשמי.
\end{abstract}
\medskip{}

\tableofcontents{}

\newpage{}
\begin{xca}
)שאלה מתרגול פעיל(

תהי \L{$f:[0,1]\to\mathbb{R}$} פונקציה רציפה כך ש-\L{$\int_{0}^{1}f(x)dx=0$}. 

הוכיחו שקיים \L{$c\in[0,1]$} כך ש-\L{$f(c)=0$}. 
\end{xca}

\begin{sol*}
נניח בשלילה כי \L{%
\[
\forall x\in\left[0,1\right]:f(x)>0
\]
}%

לפי טענה שראינו, היות והפונקציה רציפה, היא אינטגרבילית.

אז בפרט האינטגרל שלה שווה לאינטגרל התחתון שלה, שהוא הסופרמום של קבוצת
סכומי הדארבו התחתונים שלה.

נבחר חלוקה \L{$P$} ונסתכל על תת-קטע כלשהו \L{$\left[x_{i-1},x_{i}\right]$}. 

זהו תת-קטע לא מנוון וכל ערכי הפונקציה בו גדולים ממש מאפס.

בנוסף, זהו קטע סגור וחסום והפונ' רציפה, לכן לפי המשפט השני של ויירשטראס
האינפימום של ערכי הפונקציה בו שווה למינימום של ערכי הפונקציה בו, מכאן
האינפימום גם הוא גדול ממש מאפס.

נקבל \L{%
\[
m_{i}\left(x_{i}-x_{i-1}\right)>0
\]
}%

היות והפונקציה חיובית בכל הקטע, אין איבר בסכום \L{$L\left(f,P\right)$}
בעל סימן שלילי.

מכאן, \L{$L\left(f,P\right)>0$} ולכן \L{%
\[
\int_{0}^{1}f(x)\ dx=\underline{\int_{0}^{1}}f(x)\ dx\ge L\left(f,P\right)>0
\]
}%

בסתירה לנתון.

באופן דומה, נניח בשלילה כי \L{%
\[
\forall x\in\left[0,1\right]:f(x)<0
\]
}%

נגדיר את הפונקציה \L{$g(x)=-f(x)$} ונקבל לפי מה שהוכחנו \L{%
\[
\int_{0}^{1}g(x)\ dx>0
\]
}%

אז לפי הומוגניות\L{%
\[
\int_{0}^{1}f(x)\ dx=-\int_{0}^{1}g(x)\ dx<0
\]
}%

בסתירה לנתון.

משתי הסתירות נקבל שקיים \L{$x_{1}$} עבורו \L{$f\left(x_{1}\right)\le0$}
וקיים \L{$x_{2}$} עבורו \L{$f\left(x_{2}\right)\ge0$}.

לכן, לפי משפט ערך הביניים, קיים \L{$c\in\left[0,1\right]$} עבורו
\L{$f\left(c\right)=0$} כנדרש.
\end{sol*}
\begin{xca}
)שאלה מתרגול פעיל(

בכל סעיף נתונה \L{$f:[0,6]\to\mathbb{R}$}.

וודאו ש-\L{$f$} אינטגרבילית ב-\L{$[0,6]$} וקבלו ביטוי מפורש )ללא
סימן האינטגרל( לפונקציה \L{$\tilde{F}:[0,6]\to\mathbb{R}$} המוגדרת
על ידי \L{$\tilde{F}(x)=\int_{0}^{x}f(t)dt$}.

מצאו את תחום הגזירות של \L{$\tilde{F}$} וחשבו את \L{$\tilde{F}'$}.
\L{%
\[
(i)\ \forall t\in[0,6]\quad f(t)=1\qquad\qquad(ii)\ f(t)=\begin{cases}
0 & 0\le t<4\\
1 & t=4\\
0 & 4<t\le6
\end{cases}
\]
}%
\end{xca}

\begin{sol*}
\L{$\left(i\right)$} לפי תכונות האינטגרל המסוים )תכונה: יחידה(, הפונקציה
\L{$1$} אינטגרבילית ב-\L{$\left[0,6\right]$} ומתקיים\L{%
\[
\int_{0}^{x_{0}}f(x)\ dx=\int_{0}^{x_{0}}1(x)\ dx=x_{0}-0=x_{0}
\]
}%

כלומר \L{$\tilde{F}:\left[0,6\right]\rightarrow\mathbb{R}$} זו הפונקציה
\L{$\tilde{F}\left(x\right)=x$}.

היא גזירה בכל התחום )עם גזירה חד-צדדית בקצוות( ונגזרתה שווה לאחד.

\L{$\left(ii\right)$} נסתכל על הפונקציה \L{$g(t)=0$} בקטע \L{$\left[0,6\right]$}.

אז לפי תכונת ההומוגניות של האינטגרל המסוים\L{%
\[
\int_{0}^{x_{0}}0(x)\ dx=\int_{0}^{x_{0}}0\cdot1(x)\ dx=0\cdot\int_{0}^{x_{0}}1(x)\ dx=0
\]
}%

ניזכר בטענה: שינוי ערכי פונקציה אינטגרבילית במספר סופי של נקודות לא
משנה את האינטגרביליות ולא את האינטגרל.

מכאן, גם הפונקציה \L{$f(t)$} אינטגרבילית ומתקיים\L{%
\[
\int_{0}^{x_{0}}f(x)\ dx=0
\]
}%

כלומר \L{$\tilde{F}:\left[0,6\right]\rightarrow\mathbb{R}$} זו הפונקציה
\L{$\tilde{F}\left(x\right)=0$}.

היא גזירה בכל התחום )עם גזירה חד-צדדית בקצוות( ונגזרתה שווה לאפס.
\end{sol*}
\begin{xca}
)שאלה מתרגול פעיל(

יהיו \L{$a,b\in\mathbb{R}$} כך ש- \L{$a<b$} ו-\L{$f,g:[a,b]\to\mathbb{R}$}
פונקציות חסומות המקיימות \L{$f(x)\le g(x)$} לכל \L{$x\in[a,b]$}.

הוכיחו או הפריכו:

)א( \L{%
\[
\underline{\int_{a}^{b}}f(x)dx\le\underline{\int_{a}^{b}}g(x)dx
\]
}%

)ב( \L{%
\[
\overline{\int_{a}^{b}}f(x)dx\le\overline{\int_{a}^{b}}g(x)dx
\]
}%
\end{xca}

\begin{sol*}
)א( תהי \L{$P$} חלוקה של הקטע \L{$\left[a,b\right]$}.

נסתכל על \L{$L\left(f,P\right)$} ועל \L{$L\left(g,P\right)$}.

נשים לב שאנו משתמשים באותה חלוקה עבור שתי הפונ'.

נסתכל על קטע \L{$\left[x_{i-1},x_{i}\right]$} כלשהו.

מהנתון, לכל \L{$x$} בקטע מתקיים \L{$f(x)\le g(x)$}.

מהגדרת האינפימום, \L{$m_{f,i}\le f(x)\le g(x)$}.

כלומר, \L{$m_{f,i}$} הוא חסם מלרע של ערכי \L{$g(x)$} בקטע.

לכן, האינפימום של ערכי \L{$g(x)$} בקטע גדול או שווה אליו, ונקבל \L{$m_{f,i}\le m_{g,i}$}.

מכאן, לכל \L{$i$} מתקיים \L{$m_{f,i}\left(x_{i}-x_{i-1}\right)\le m_{g,i}\left(x_{i}-x_{i-1}\right)$}.

נסכום את אי-השוויונות ונקבל\L{%
\[
\underset{i=1}{\overset{n}{\sum}}m_{f,i}\left(x_{i}-x_{i-1}\right)\le\underset{i=1}{\overset{n}{\sum}}m_{g,i}\left(x_{i}-x_{i-1}\right)
\]
}%

כלומר\L{%
\[
L\left(f,P\right)\le L\left(g,P\right)
\]
}%

האינטגרל התחתון הוא הסופרמום של קבוצת סכומי הדרבו התחתונים, מכאן לכל
חלוקה \L{$P$} מתקיים\L{%
\[
L\left(f,P\right)\le L\left(g,P\right)\le\underline{\int_{a}^{b}}g(x)dx
\]
}%

קיבלנו ש-\L{$\underline{\int_{a}^{b}}g(x)dx$} הוא חסם מלעיל של הקבוצה
\L{$\left\{ L\left(f,P\right)\text{ s.t P is a partition}\right\} $}.

לכן, מהגדרת הסופרמום, הסופרמום של הקבוצה קטן או שווה אליו. 

בסך הכל, זה יוצא בדיוק\L{%
\[
\underline{\int_{a}^{b}}f(x)dx\le\underline{\int_{a}^{b}}g(x)dx
\]
}%

כנדרש.

)ב( תהי \L{$P$} חלוקה של הקטע \L{$\left[a,b\right]$}.

נסתכל על \L{$U\left(f,P\right)$} ועל \L{$U\left(g,P\right)$}.

אנו משתמשים באותה חלוקה, לכן אורכי תת הקטעים שווים.

נסתכל על קטע \L{$\left[x_{i-1},x_{i}\right]$} כלשהו.

מהנתון, לכל \L{$x$} בקטע מתקיים \L{$f(x)\le g(x)$}.

מהגדרת הסופרמום, \L{$f(x)\le g(x)\le M_{g,i}$}.

כלומר, \L{$M_{g,i}$} הוא חסם מלעיל של ערכי \L{$f(x)$} בקטע.

לכן, הסופרמום של ערכי \L{$f(x)$} קטן או שווה אליו, ונקבל \L{$M_{f,i}\le M_{g,i}$}.

מכאן, לכל \L{$i$} מתקיים \L{$M_{f,i}\left(x_{i}-x_{i-1}\right)\le M_{g,i}\left(x_{i}-x_{i-1}\right)$}.

נסכום את אי-השוויונות ונקבל\L{%
\[
\underset{i=1}{\overset{n}{\sum}}M_{f,i}\left(x_{i}-x_{i-1}\right)\le\underset{i=1}{\overset{n}{\sum}}M_{g,i}\left(x_{i}-x_{i-1}\right)
\]
}%

כלומר\L{%
\[
U\left(f,P\right)\le U\left(g,P\right)
\]
}%

האינטגרל העליון הוא האינפימום של קבוצת סכומי הדרבו העליונים, מכאן
לכל חלוקה \L{$P$} מתקיים\L{%
\[
\overline{\int_{a}^{b}}f(x)dx\le U\left(f,P\right)\le U\left(g,P\right)
\]
}%

קיבלנו ש-\L{$\overline{\int_{a}^{b}}f(x)dx$} הוא חסם מלרע של הקבוצה
\L{$\left\{ U\left(g,P\right)\text{ s.t P is a partition}\right\} $}.

לכן, מהגדרת האינפימום, האינפימום שלה גדול או שווה אליו.

בסך הכל, זה יוצא בדיוק\L{%
\[
\overline{\int_{a}^{b}}f(x)dx\le\overline{\int_{a}^{b}}g(x)dx
\]
}%

כנדרש.

)ב( הוכחה נוספת, מעט יותר מסורבלת ובשלילה, לצורך העשרה )ציטוט של המתרגל
ומפעיל קנדה-עליון-וירטואלי האגדי: ג'נטלמן לא מוכיח בשלילה כשאפשר שלא(:

הטענה שקולה למעשה לטענה\L{%
\[
I_{f}\overset{\text{def}}{=}\inf\left\{ U\left(f,P\right)\text{ s.t P is a partition}\right\} \le\inf\left\{ U\left(g,P\right)\text{ s.t P is a partition}\right\} \overset{\text{def}}{=}I_{g}
\]
}%
נניח בשלילה כי\L{%
\[
I_{f}>I_{g}
\]
}%

נבחר\L{%
\[
\varepsilon=I_{f}-I_{g}>0
\]
}%

ולפי הגדרת האינפימום קיימת חלוקה \L{$P'$} עבורה \L{%
\[
I_{g}\le U\left(g,P'\right)<I_{g}+\varepsilon=I_{f}
\]
}%

היות ו-\L{$I_{f}$} הוא האינפימום, בפרט מתקיים\L{%
\[
I_{f}\le U\left(f,P'\right)
\]
}%

כלומר, כאשר תת-הקטעים הם אותם תת-הקטעים של \L{$\left[a,b\right]$}
לפי החלוקה \L{$P'$}, קיבלנו ש-\L{%
\[
\underset{i=1}{\overset{n}{\sum}}M_{g,i}\left(x_{i}-x_{i-1}\right)<\underset{i=1}{\overset{n}{\sum}}M_{f,i}\left(x_{i}-x_{i-1}\right)
\]
}%

נשים לב שאנו משתמשים באותה חלוקה עבור שתי הפונ'.

נסתכל על קטע \L{$\left[x_{i-1},x_{i}\right]$} כלשהו.

מהנתון, לכל \L{$x$} בקטע מתקיים \L{$f(x)\le g(x)$}.

מהגדרת הסופרמום, \L{$f(x)\le g(x)\le M_{g,i}$}.

כלומר, \L{$M_{g,i}$} הוא חסם מלעיל של ערכי \L{$f(x)$} בקטע.

לכן, הסופרמום של ערכי \L{$f(x)$} קטן או שווה אליו, ונקבל \L{$M_{f,i}\le M_{g,i}$}.

מכאן, לכל \L{$i$} מתקיים \L{$M_{f,i}\left(x_{i}-x_{i-1}\right)\le M_{g,i}\left(x_{i}-x_{i-1}\right)$}.

מאיחוד אי-השוויונות נקבל את הסתירה\L{%
\[
\underset{i=1}{\overset{n}{\sum}}M_{f,i}\left(x_{i}-x_{i-1}\right)\le\underset{i=1}{\overset{n}{\sum}}M_{g,i}\left(x_{i}-x_{i-1}\right)
\]
}%
\end{sol*}
\begin{xca}
)שאלה מתרגול פעיל(

יהיו \L{$\alpha,\beta\in\mathbb{R}$} ויהי \L{$D\subseteq\mathbb{R}$}
מקטע.

הפריכו או הוכיחו את הטענות הבאות:

)א( אם הפונקציות \L{$f$} ו-\L{$g$} רציפות במידה שווה ב-\L{$D$},
אז הפונקציה \L{$\alpha f+\beta g$} רציפה במידה שווה ב-\L{$D$}.

)ב( אם פונקציה \L{$f$} רציפה במידה שווה ב-\L{$D$}, אז הפונקציה \L{$|f|$}
גם היא רציפה במידה שווה ב-\L{$D$}. 
\end{xca}

\begin{sol*}
)א( יהי \L{$f$} ו-\L{$g$} פונ' רציפות במידה שווה. 

יהי \L{$\varepsilon>0$}. נרצה להוכיח כי קיימת \L{$\delta>0$} עבורה\L{%
\[
\forall x_{1},x_{2}\in D:\left|x_{1}-x_{2}\right|<\delta\Rightarrow\left|\left(\alpha f+\beta g\right)\left(x_{1}\right)-\left(\alpha f+\beta g\right)\left(x_{2}\right)\right|<\varepsilon
\]
}%
הביטוי הנגרר שקול לביטוי\L{%
\[
\left|\alpha\cdot f\left(x_{1}\right)-\alpha\cdot f\left(x_{2}\right)+\beta\cdot g\left(x_{1}\right)-\beta\cdot g\left(x_{2}\right)\right|<\varepsilon
\]
}%

לפי א\textquotedblright ש המשולש\L{%
\[
\left|\alpha\cdot f\left(x_{1}\right)-\alpha\cdot f\left(x_{2}\right)+\beta\cdot g\left(x_{1}\right)-\beta\cdot g\left(x_{2}\right)\right|\le\left|\alpha\cdot f\left(x_{1}\right)-\alpha\cdot f\left(x_{2}\right)\right|+\left|\beta\cdot g\left(x_{1}\right)-\beta\cdot g\left(x_{2}\right)\right|
\]
}%
לכן מספיק להוכיח שנובע \L{%
\[
\left|\alpha\cdot f\left(x_{1}\right)-\alpha\cdot f\left(x_{2}\right)\right|+\left|\beta\cdot g\left(x_{1}\right)-\beta\cdot g\left(x_{2}\right)\right|<\varepsilon
\]
}%
שזה שקול לביטוי\L{%
\[
\left|\alpha\right|\left|f\left(x_{1}\right)-f\left(x_{2}\right)\right|+\left|\beta\right|\left|g\left(x_{1}\right)-g\left(x_{2}\right)\right|<\varepsilon
\]
}%

אם \L{$\left|\alpha\right|=0$} אז \L{$\alpha=0$} והפונקציה \L{$\alpha f$}
היא פונ' האפס לכן נותר להוכיח רק עבור \L{$\beta g$}.

במקרה זה אפשר להשתמש רק בחצי השני של ההוכחה שמניחה \L{$\left|\alpha\right|,\left|\beta\right|>0$}.

אם \L{$\left|\beta\right|=0$} אז \L{$\beta=0$} והפונקציה \L{$\beta g$}
היא פונ' האפס לכן נותר להוכיח רק עבור \L{$\alpha f$}.

במקרה זה אפשר להשתמש רק בחצי הראשון של ההוכחה שמניחה \L{$\left|\alpha\right|,\left|\beta\right|>0$}.

אם שתיהן פונ' האפס אז הפונקציה \L{$\alpha f+\beta g$} קבועה ולכן
רציפה במ\textquotedblright ש.

נניח כי \L{$\left|\alpha\right|,\left|\beta\right|>0$}.

אז \L{$\frac{\varepsilon}{2\left|\alpha\right|}>0$} וגם \L{$\frac{\varepsilon}{2\left|\beta\right|}>0$}.

מכך ש-\L{$f$} ו-\L{$g$} רציפות במ\textquotedblright ש, קיימות \L{$\delta_{1},\delta_{2}>0$}
כך ש-\L{%
\[
\forall x_{1},x_{2}\in D:\left|x_{1}-x_{2}\right|<\delta_{1}\Rightarrow\left|f\left(x_{1}\right)-f\left(x_{2}\right)\right|<\frac{\varepsilon}{2\left|\alpha\right|}
\]
}%
\L{%
\[
\forall x_{1},x_{2}\in D:\left|x_{1}-x_{2}\right|<\delta_{2}\Rightarrow\left|g\left(x_{1}\right)-g\left(x_{2}\right)\right|<\frac{\varepsilon}{2\left|\beta\right|}
\]
}%

מכאן\L{%
\[
\forall x_{1},x_{2}\in D:\left|x_{1}-x_{2}\right|<\delta_{1}\Rightarrow\left|\alpha\right|\left|f\left(x_{1}\right)-f\left(x_{2}\right)\right|<\frac{\varepsilon}{2}
\]
}%
\L{%
\[
\forall x_{1},x_{2}\in D:\left|x_{1}-x_{2}\right|<\delta_{2}\Rightarrow\left|\beta\right|\left|g\left(x_{1}\right)-g\left(x_{2}\right)\right|<\frac{\varepsilon}{2}
\]
}%

נבחר \L{$\delta=\min\left\{ \delta_{1},\delta_{2}\right\} $}.

יהיו \L{$x_{1},x_{2}\in D$} כך ש- \L{$\left|x_{1}-x_{2}\right|<\delta$}.

אז \L{%
\[
\left|\alpha\right|\left|f\left(x_{1}\right)-f\left(x_{2}\right)\right|+\left|\beta\right|\left|g\left(x_{1}\right)-g\left(x_{2}\right)\right|<\frac{\varepsilon}{2}+\frac{\varepsilon}{2}=\varepsilon
\]
}%

כנדרש.

)ב( יהי \L{$\varepsilon>0$}. הפונ' \L{$f$} רציפה במ\textquotedblright ש
לכן \L{$\exists\delta>0$} עבורה לכל \L{$x_{1},x_{2}\in D$} מתקיים\L{%
\[
\left|x_{1}-x_{2}\right|<\delta\Rightarrow\left|f\left(x_{1}\right)-f\left(x_{2}\right)\right|<\varepsilon
\]
}%

נניח כי \L{$\left|x_{1}-x_{2}\right|<\delta$}. אז לפי א\textquotedblright ש
המשולש ההפוך\L{%
\[
\left|\left|f\left(x_{1}\right)\right|-\left|f\left(x_{2}\right)\right|\right|\le\left|f\left(x_{1}\right)-f\left(x_{2}\right)\right|<\varepsilon
\]
}%

כנדרש.
\end{sol*}
\begin{xca}
תהא \L{$f:\left[a,b\right]\rightarrow\mathbb{R}$} פונקציה מונוטונית. 

הוכיחו כי \L{$f$} אינטגרבילית לפי רימן.
\end{xca}

\begin{proof}
כרגע לא בעמוד.
\end{proof}
\begin{xca}
תהא \L{$f:\left[a,b\right]\rightarrow\mathbb{R}$} פונקציה רציפה. 

הוכיחו כי \L{$f$} אינטגרבילית לפי רימן.
\end{xca}

\begin{proof}
כרגע לא בעמוד.
\end{proof}
\begin{xca}
)שאלה מתרגול פעיל(

)א( חשבו את האינטגרל \L{$\int_{0}^{\pi}\sin^{2}(x)dx$} בשתי דרכים:

)i( ישירות, באמצעות הזהות \L{$\cos(2x)=\cos^{2}(x)-\sin^{2}(x)=1-2\sin^{2}(x)$}.

)ii( באמצעות אינטגרציה בחלקים ושימוש בזהות \L{$\cos^{2}(x)+\sin^{2}(x)=1$}.

)ב( מהו הערך הממוצע של \L{$f(x)=\sin^{2}(x)$} בקטע \L{$[-\pi,\pi]$}? 
\end{xca}

\begin{sol*}
כרגע לא בעמוד.
\end{sol*}
\begin{xca}
)שאלה מתרגול פעיל(

יהיו \L{$a,b\in\mathbb{R}$} כך ש- \L{$a<b$} ותהי \L{$f:[a,b]\to\mathbb{R}$}
פונקציה רציפה.

נתון כי \L{$\int_{a}^{b}f(x)g(x)dx=0$} עבור כל פונקציה אינטגרבילית
\L{$g:[a,b]\to\mathbb{R}$}.

הוכיחו כי \L{$f(x)=0$} לכל \L{$x\in[a,b]$}.
\end{xca}

\begin{proof}
כרגע לא בעמוד.
\end{proof}
\begin{xca}
)שאלה מתרגול פעיל(

)א( הוכיחו \textbackslash{} הפריכו: אם \L{$f:[1,\infty)\to\mathbb{R}$}
אינטגרבילית ב-\L{$[a,b]$} עבור \uline{כל} \L{$a,b\in\mathbb{R}$}
המקיימים \L{$1\leq a<b$}, אזי \L{$f$} חסומה.

)ב( הוכיחו \textbackslash{} הפריכו: אם \L{$f$} ו-\L{$g$} רציפות
במידה שווה ב-\L{$\mathbb{R}$}, אז גם פונקציית ההרכבה \L{$g\circ f$}
רציפה במידה שווה ב-\L{$\mathbb{R}$}.

)ג( הוכיחו \textbackslash{} הפריכו: אם \L{$f$} פונקציה חיובית )דהיינו
\L{$0<f(x)$} לכל \L{$x$} בקטע \L{$[a,b]$}(, אז \L{$\ln(f)$} אינטגרבילית
בקטע \L{$[a,b]$}. 
\end{xca}

\begin{sol*}
כרגע לא בעמוד.
\end{sol*}
\begin{xca}
)שאלה מתרגול פעיל(

חשבו את ערכו של הגבול \L{%
\[
\lim_{x\to1}\frac{\int_{1}^{x}e^{t^{2}}dt}{\sin(x-1)}
\]
}%
)אל תנסו למצוא ביטוי מפורש ללא סימן האינטגרל עבור המונה(
\end{xca}

\begin{sol*}
כרגע לא בעמוד.
\end{sol*}
\begin{xca}
הוכיחו את המשפט היסודי של החדו\textquotedblright א.
\end{xca}

\begin{proof}
כרגע לא בעמוד.
\end{proof}
\begin{xca}
הוכיחו כי פונקציית רימן רציפה בכל נקודה אי-רציונלית.
\end{xca}

\begin{proof}
כרגע לא בעמוד.
\end{proof}

\end{document}
